\documentclass[11pt]{article}
\usepackage[margin=1in]{geometry}
\usepackage{amsmath,amssymb,amsthm,booktabs}
\usepackage[T1]{fontenc}
\usepackage{lmodern}
\usepackage{microtype}
\usepackage[title]{appendix}
\DeclareUnicodeCharacter{2227}{\ensuremath{\land}}
\DeclareUnicodeCharacter{00AC}{\ensuremath{\neg}}
\DeclareUnicodeCharacter{211D}{\ensuremath{\mathbb{R}}}
\DeclareUnicodeCharacter{2192}{\ensuremath{\rightarrow}}
\DeclareUnicodeCharacter{2200}{\ensuremath{\forall}}
\DeclareUnicodeCharacter{2208}{\ensuremath{\in}}
\DeclareUnicodeCharacter{2264}{\ensuremath{\leq}}
\DeclareUnicodeCharacter{00D7}{\ensuremath{\times}}
\DeclareUnicodeCharacter{02E2}{\ensuremath{{}^{\mathrm{s}}}}
\DeclareUnicodeCharacter{03BC}{\ensuremath{\mu}}
\DeclareUnicodeCharacter{2203}{\ensuremath{\exists}}
\DeclareUnicodeCharacter{207B}{\ensuremath{{}^{-}}}
\DeclareUnicodeCharacter{00B9}{\ensuremath{{}^{1}}}
\DeclareUnicodeCharacter{2260}{\ensuremath{\ne}}
\DeclareUnicodeCharacter{22A4}{\ensuremath{\top}}
\usepackage{hyperref}
\hypersetup{colorlinks=true,linkcolor=blue,citecolor=blue,urlcolor=blue,
pdfauthor={P. Graham},
pdftitle={A counterexample on two kinds of positive dependence}}
\newtheorem{theorem}{Theorem}
\newtheorem{lemma}[theorem]{Lemma}
\title{A counterexample on two kinds of positive dependence}
\author{P. Graham\\ \small\url{https://orcid.org/0009-0007-8493-1505}}
\date{October 31, 2026}
\begin{document}
\maketitle
\begin{abstract}
We construct a strict bivariate Archimedean copula that is positively
quadrant dependent but is not positive measure inducing. This disproves
the forward implication in the equivalence conjectured by Sebastian
Fuchs and Marco Tschimpke. Their work introduced the positive measure
inducing property through a reflected copula sum and established the
equivalence for the Frank family. Our example separates the two
properties in the full Archimedean class. Its inverse generator is the
maximum of two exponential functions. Convexity gives copula validity,
and a piecewise linear inequality in logarithmic coordinates gives
positive quadrant dependence throughout the unit square. A rectangle
strictly inside the open square $(0,1/2)^2$ has a negative increment for the
reflected sum, excluding the positive measure required by the definition.
All corner estimates use rational arithmetic and squaring, without a
density assumption. The complete argument, including generator
admissibility, boundary values and the measure-theoretic contradiction,
is formalized in Lean. The converse implication and the separate
negative-dependence equivalence are not resolved by this example.

This work has not yet been reviewed by experts.
The proofs, the Lean formalization and the text were produced by AI models
(from OpenAI and Anthropic), working under the author's direction.
\end{abstract}
\medskip
\noindent\textit{Mathematics Subject Classification (2020):} 60E05, 62H05.\\
\textit{Keywords:} Archimedean copula, positive quadrant dependence,
positive measure inducing dependence.

\section{Introduction}
\begin{theorem}\label{thm:main}
For $t\geq0$ and $0<u\leq1$, put
\[
 \psi(t)=\max\{e^{-t},\tfrac14 e^{-t/2}\},\qquad
 \phi(u)=\max\{-\log u,-2\log u-\log16\}.
\]
The function $\phi$, extended by $\phi(0)=\infty$, is a strict Archimedean
generator in the sense of \cite{ft}, and $\psi$ is its inverse on
$[0,\infty)$. The function on $[0,1]^2$ defined by
\[
 C(u,v)=
 \begin{cases}\psi(\phi(u)+\phi(v)),&u,v>0,\\0,&uv=0\end{cases}
\]
is a copula, satisfies $C(u,v)\geq uv$ everywhere, and is not positive
measure inducing in the sense of Fuchs and Tschimpke.
\end{theorem}

A copula is a two-increasing function on the unit square, zero on the
lower edges and with uniform margins $C(u,1)=u$ and $C(1,v)=v$.
The inequality in the theorem is positive quadrant dependence (PQD).
For example, the independence copula $\Pi(u,v)=uv$ has PQD with equality.
For any function $F$ of two variables write its rectangle increment as
\[
 \Delta F[a,b;d,e]=F(b,e)-F(a,e)-F(b,d)+F(a,d).
\]
Two-increasing means that these increments are nonnegative whenever
$a\leq b$ and $d\leq e$ in the unit interval.

Fuchs and Tschimpke's Definition 1 and equation (2.2) in
\cite{ft}, with numbering from the arXiv version specified there,
define the positive measure inducing property (PMI) by requiring that
\begin{equation}\label{eq:E}
 E_C(u,v)=C(u,v)+C(1-u,v)+C(u,1-v)+C(1-u,1-v)-1
\end{equation}
induce a positive measure on $(0,1/2)^2$. Thus every half-open rectangle
$(a,b]\times(d,e]$ compactly contained in that square must have
nonnegative $E_C$-increment. For independence, $E_\Pi$ is identically zero.
Example 6 of \cite{ft} gives the PQD--PMI equivalence for Frank copulas.
Conjecture 1 proposes it for all Archimedean copulas; Theorem~\ref{thm:main}
disproves the implication PQD $\Rightarrow$ PMI. The conjecture's reverse
implication and its separate NMI equivalence are outside this conclusion.

The generator description is classical: the bivariate case of
McNeil and Ne\v{s}lehov\'a's Theorem 2.2 in \cite{mn} uses decreasing convex
inverse generators. We verify those properties and the copula axioms
directly. The distinction between a positive measure and a density is
essential: Remark 1(2) of \cite{ft} gives a density criterion under absolute
continuity, whereas our proof uses Definition 1 itself. Lean checks the
full statement from these definitions, as described in Section~\ref{sec:formal}.
The superadditivity criterion behind our PQD argument appears in
M\"uller and Scarsini's Theorem 2.7 in the working version of \cite{ms},
attributed there to Joe, \emph{Multivariate Models and Dependence Concepts}
(1997), page 109.
Fern\'andez S\'anchez and Trutschnig's Corollary 4 in the preprint of
\cite{fst} expresses the mass of an Archimedean level curve in terms of
one-sided generator derivatives. That formula motivates attention to
corners; here direct rectangle increments suffice.

\section{Proof overview}
The maximum of two exponentials remains convex and strictly decreasing.
In logarithmic coordinates its inverse has two affine branches, whose
maximum is superadditive. These two observations establish the global
copula and PQD assertions. The separation of PMI occurs at the change of
branch: near a reflected corner, the copula is a maximum of a product
and a square root of that product. The alternating sum of sixteen corner
values over one small rectangle is negative. Exact rational enclosures
make this last sign calculation elementary and avoid differentiating
across the change of branch.

\section{The generator and positive quadrant dependence}
\begin{lemma}\label{lem:gen}
On $(0,1]$ and $[0,\infty)$, respectively, the functions $\phi$ and $\psi$
in Theorem~\ref{thm:main} are continuous, strictly decreasing and convex. They are inverse to each other,
with $\phi(1)=0$, $\phi(u)\to\infty$ as $u\downarrow0$,
$\psi(0)=1$ and $\psi(t)\to0$ as $t\to\infty$.
Their construction defines a copula on the closed unit square.
\end{lemma}
\begin{proof}
Each branch has the claimed continuity, monotonicity and convexity;
taking a maximum preserves these properties, including strict decrease.
The branches of $\psi$ meet at $t=\log16$, at value $1/16$.
Inverting on the two sides gives $-\log u$ for $u\geq1/16$ and
$-2\log u-\log16$ for $u\leq1/16$, exactly $\phi$.
The endpoint values and limits follow from these formulas.

The inverse identities give the margins, and the definition gives the
zero edges. For a positive rectangle with $a\leq b$ and $d\leq e$ set
\[
 s=\phi(b)+\phi(e),\quad h=\phi(a)-\phi(b)\geq0,\quad
 k=\phi(d)-\phi(e)\geq0.
\]
Its $C$-increment is
$\psi(s+h+k)-\psi(s+h)-\psi(s+k)+\psi(s)$, which is nonnegative:
for a convex function, an increment of fixed nonnegative length is
nondecreasing in its starting point. If $a=0$, the increment is
$C(b,e)-C(b,d)\geq0$ by monotonicity; the case $d=0$ is the same,
and repeated endpoints give zero. Thus all closed-square rectangle
increments are nonnegative. The margins and groundedness also imply
$0\leq C(u,v)\leq\min(u,v)$.
\end{proof}

\begin{lemma}\label{lem:pqd}
The copula in Theorem~\ref{thm:main} satisfies $C(u,v)\geq uv$ for every
$(u,v)\in[0,1]^2$.
\end{lemma}
\begin{proof}
Let $L=\log4$ and $f(x)=\max(x,2x-2L)$ for $x\geq0$.
This function is superadditive. If $x,y\leq2L$, then
$f(x)+f(y)=x+y\leq f(x+y)$. If $x\geq2L\geq y$, then
$f(x)+f(y)=2x-2L+y\leq2(x+y)-2L\leq f(x+y)$.
The symmetric case is identical. If both exceed $2L$, the sum is
$2(x+y)-4L\leq f(x+y)$.
For positive $u,v\leq1$, take $x=-\log u$ and $y=-\log v$.
Then $\phi(u)+\phi(v)\leq\phi(uv)$, so the decreasing inverse gives
$C(u,v)\geq\psi(\phi(uv))=uv$. The zero edges give equality.
\end{proof}

\section{A negative reflected rectangle}
Set
\[
 a=\frac{63}{256},\quad b=\frac{65}{256},\quad
 d=\frac{61}{768},\quad e=\frac{67}{768}.
\]
These endpoints satisfy $0<a<b<1/2$ and $0<d<e<1/2$.
Every coordinate of each reflected corner lies in $[1/16,1]$.
On this interval $\phi(u)=-\log u$, hence
\begin{equation}\label{eq:corner}
 C(u,v)=\max\{uv,\sqrt{uv}/4\}.
\end{equation}
Put $D=196608$ and $q(n)=\sqrt{n/D}/4$. Table~\ref{tab:corners}
lists all sixteen values needed in \eqref{eq:E}. The branch follows by
comparing $uv$ with $1/16$, or its numerator with $D/16=12288$.

\begin{table}[ht]
\centering
\begin{tabular}{ccccc}
\toprule
$(u,v)$ & $C(u,v)$ & $C(1-u,v)$ & $C(u,1-v)$ & $C(1-u,1-v)$\\
\midrule
$(b,e)$ & $q(4355)$ & $12797/D$ & $45565/D$ & $133891/D$\\
$(a,e)$ & $q(4221)$ & $12931/D$ & $44163/D$ & $135293/D$\\
$(b,d)$ & $q(3965)$ & $q(11651)$ & $45955/D$ & $135037/D$\\
$(a,d)$ & $q(3843)$ & $q(11773)$ & $44541/D$ & $136451/D$\\
\bottomrule
\end{tabular}
\caption{Reflected corner values for the negative rectangle.}\label{tab:corners}
\end{table}

\begin{lemma}\label{lem:neg}
For the displayed rectangle, $\Delta E_C[a,b;d,e]<0$.
\end{lemma}
\begin{proof}
The following integers $k$ satisfy $k/10^6\leq q(n)\leq(k+1)/10^6$:
\[
\begin{array}{c|rrrrrr}
n&4355&4221&3965&11651&3843&11773\\\hline
k&37207&36630&35502&60858&34952&61176
\end{array}
\]
Each enclosure follows by squaring nonnegative numbers: its two checks
are $16Dk^2\leq10^{12}n\leq16D(k+1)^2$.
In the alternating sum of Table~\ref{tab:corners}, use upper endpoints
in the first and fourth rows, and lower endpoints in the second and
third rows. The constant terms in \eqref{eq:E} cancel, and this gives
\[
 \Delta E_C[a,b;d,e]\leq
 \frac{37208-36630-35502-60858+34953+61177}{10^6}
 -\frac{134}{196608}
 =-\frac{512347}{1536000000}<0.
\]
\end{proof}

\begin{proof}[Proof of Theorem~\ref{thm:main}]
Lemma~\ref{lem:gen} establishes strict generator admissibility and all
copula axioms, and Lemma~\ref{lem:pqd} proves PQD. If $C$ were PMI, its
induced positive measure on the open square $(0,1/2)^2$ would assign the
rectangle $(a,b]\times(d,e]$ the value $\Delta E_C[a,b;d,e]$.
Lemma~\ref{lem:neg} makes this value negative, a contradiction.
\end{proof}

\section{Formal verification}\label{sec:formal}
The complete theorem has the following Lean statement, in namespace
\texttt{FT23}:
\begin{verbatim}
theorem conjecture_counterexample :
    StrictGenerator phi psi ∧ IsCopula C ∧
    PQD C ∧ ¬ PrimaryPMI C
\end{verbatim}
Here the logical symbols are conjunction and negation.
\texttt{StrictGenerator} records the continuity, strict decrease,
convexity, inverse identities and endpoint limits of both displayed
functions. \texttt{IsCopula} includes groundedness, uniform margins and
two-increasingness for all rectangles in the closed square, including
repeated endpoints. \texttt{PQD} quantifies over that whole square.
\texttt{PrimaryPMI} asserts the existence of a positive Borel measure on
the open square $(0,1/2)^2$ whose finite half-open rectangle values equal
the corresponding increments of $E_C$. Thus its negation has the
measure-theoretic meaning used above, rather than merely rejecting a
smooth density test. Appendix~\ref{app:definitions} quotes the definitions
and makes the endpoint convention explicit.

The formal proof establishes the generator identities and analytic
inequalities, then bounds the same sixteen corner values by exact
rational numbers. Its enclosures are tighter than the displayed ones.
The independence example checks that $E_\Pi$ and all its increments
vanish.
Only Lean's standard axioms of propositional extensionality, classical
choice and quotient soundness occur. The build checks this axiom list
and the final theorem's type; no compiled numerical evaluation is used.

The accompanying Lean project contains all proof sources and pinned
library dependencies. Its README gives the standard Lake commands;
after the dependencies are available, checking takes a few minutes.
Reproducing the theorem requires Lean and mathlib but no search code,
external certificate, or unproved analytic premise.

\begin{thebibliography}{9}
\bibitem{ft}
S. Fuchs and M. Tschimpke,
\emph{A novel positive dependence property and its impact on a popular
class of concordance measures},
Journal of Multivariate Analysis \textbf{200} (2024), 105259.
\href{https://doi.org/10.1016/j.jmva.2023.105259}{doi:10.1016/j.jmva.2023.105259}.
Statement numbering follows
\href{https://arxiv.org/abs/2306.09676v1}{arXiv:2306.09676v1}, June 16, 2023.

\bibitem{mn}
A. J. McNeil and J. Ne\v{s}lehov\'a,
\emph{Multivariate Archimedean copulas, $d$-monotone functions and
$\ell_1$-norm symmetric distributions},
Annals of Statistics \textbf{37} (2009), 3059--3097.
\href{https://doi.org/10.1214/07-AOS556}{doi:10.1214/07-AOS556}.
Theorem 2.2; \href{https://arxiv.org/abs/0908.3750v1}{arXiv:0908.3750v1}.
\bibitem{ms}
A. M\"uller and M. Scarsini,
\emph{Archimedean copulae and positive dependence},
Journal of Multivariate Analysis \textbf{93} (2005), 434--445.
\href{https://doi.org/10.1016/j.jmva.2004.04.003}{doi:10.1016/j.jmva.2004.04.003}.
Theorem 2.7 in the
\href{https://www.bemservizi.unito.it/repec/icr/wp2003/Scarsini25-03.pdf}{2003 working version}.

\bibitem{fst}
J. Fern\'andez S\'anchez and W. Trutschnig,
\emph{Singularity aspects of Archimedean copulas},
Journal of Mathematical Analysis and Applications \textbf{432} (2015), 103--113.
\href{https://doi.org/10.1016/j.jmaa.2015.06.036}{doi:10.1016/j.jmaa.2015.06.036}.
Corollary 4 in the \href{https://www.trutschnig.net/p21.pdf}{June 22, 2015 preprint}.
\end{thebibliography}

\begin{appendices}
\section{Lean definitions}\label{app:definitions}
The following declarations are quoted verbatim from the formalization,
in namespace \texttt{FT23}; \texttt{Set}, \texttt{Filter} and
\texttt{MeasureTheory} are open. The real-valued Lean \texttt{phi} is used
only on $(0,1]$: since \texttt{Real.log 0 = 0}, its value at zero is zero,
while the extended value $\phi(0)=\infty$ is represented by the
\texttt{StrictGenerator.phi\_at\_zero} right-limit condition; the zero
edges of \texttt{C} are defined directly.

The definition of \texttt{C} is the strict Archimedean construction of
\cite[Section 2]{ft}, with those zero edges.
\begingroup
\footnotesize
\begin{verbatim}
def C (u v : ℝ) : ℝ := if 0 < u ∧ 0 < v then psi (phi u + phi v) else 0

def delta (F : ℝ → ℝ → ℝ) (a b d e : ℝ) : ℝ :=
  F b e - F a e - F b d + F a d
\end{verbatim}
\endgroup
Here \texttt{delta} is the half-open rectangle increment used to
interpret the induced measure in \cite[Definition 1]{ft}.
The reflection definitions below give \texttt{E}; the proved identity
\texttt{E\_expanded} is exactly \cite[equation (2.2)]{ft}.
\begingroup
\footnotesize
\begin{verbatim}
def nu1 (F : ℝ → ℝ → ℝ) (u v : ℝ) : ℝ := v - F (1-u) v
def nu2 (F : ℝ → ℝ → ℝ) (u v : ℝ) : ℝ := u - F u (1-v)
def nu (F : ℝ → ℝ → ℝ) (u v : ℝ) : ℝ := u + v - 1 + F (1-u) (1-v)
def E (F : ℝ → ℝ → ℝ) (u v : ℝ) : ℝ :=
  F u v - nu1 F u v - nu2 F u v + nu F u v

lemma E_expanded (F : ℝ → ℝ → ℝ) (u v : ℝ) :
    E F u v = F u v + F (1-u) v + F u (1-v) + F (1-u) (1-v) - 1 := by
  unfold E nu1 nu2 nu
  ring
\end{verbatim}
\endgroup
The definition of \texttt{PQD} is the inequality in
\cite[Section 2]{ft}, extended to the closed square, where the margins
and zero edges give the boundary cases.
\begingroup
\footnotesize
\begin{verbatim}
def PQD (F : ℝ → ℝ → ℝ) : Prop :=
  ∀ u ∈ Icc (0:ℝ) 1, ∀ v ∈ Icc (0:ℝ) 1, u*v ≤ F u v
\end{verbatim}
\endgroup
The predicate \texttt{ValidRect} places all four endpoints strictly
inside the square of \cite[Definition 1]{ft}; \texttt{HalfSquare} is that
square as a subtype, and \texttt{rect} is a half-open rectangle.
\begingroup
\footnotesize
\begin{verbatim}
def ValidRect (a b d e : ℝ) : Prop :=
  0 < a ∧ a < b ∧ b < 1/2 ∧ 0 < d ∧ d < e ∧ e < 1/2

def halfSquare : Set (ℝ × ℝ) := Ioo (0:ℝ) (1/2) ×ˢ Ioo (0:ℝ) (1/2)
abbrev HalfSquare := {z : ℝ × ℝ // z ∈ halfSquare}
def rect (a b d e : ℝ) : Set (ℝ × ℝ) := Ioc a b ×ˢ Ioc d e
\end{verbatim}
\endgroup
The predicate \texttt{PrimaryPMI} is the necessary rectangle-value
consequence of \cite[Definition 1]{ft}: a positive Borel measure on the
open square has each indicated finite value equal to the corresponding
increment of \texttt{E}. The clause \texttt{≠ ⊤} excludes infinite measure
before \texttt{toReal} converts its value to a real number; negating this
necessary condition therefore rules out PMI itself.
\begingroup
\footnotesize
\begin{verbatim}
def PrimaryPMI (F : ℝ → ℝ → ℝ) : Prop :=
  ∃ μ : Measure HalfSquare, ∀ a b d e : ℝ, ValidRect a b d e →
    μ (Subtype.val ⁻¹' rect a b d e) ≠ ⊤ ∧
    (μ (Subtype.val ⁻¹' rect a b d e)).toReal = delta (E F) a b d e
\end{verbatim}
\endgroup
\end{appendices}
\end{document}
